% !TeX root = ../../Book2.tex
\chapter{非线性椭圆方程与低正则理论}
\label{b2:app:H}

\begin{chapterabstract}
可测系数通常不能逐项求导，非线性通量的线性化也可能在零梯度处退化。弱形式仍保留一个可用的结构：把解的截断或适当幂函数作为测试，通量的正性就转化为能量控制。本附录先补齐线性齐次方程从 Moser 上界、对数能量到 Harnack 比较的证明，再建立一般 \(p\) 次增长通量的截断估计与局部有界性。在具有凸势的加权 \(p\)-Laplace 模型中，直接法、单调性和序结构分别给出存在、唯一、比较与稳定性。最后用径向解和点源说明穿孔域方程、分布通量与自然能量解的区别。
\end{chapterabstract}

\begin{learninggoals}
\begin{itemize}
\item\label{b2:item:H2-goal-test}说明截断、幂函数和倒数测试的合法性，分清当前可积性与本轮估计新得到的可积性。
\item\label{b2:item:H2-goal-iteration}串联局部上界、对数平均振荡与正负小幂，证明线性 Harnack 比较，并辨认它与 Hölder 正则性的接口。
\item\label{b2:item:H2-goal-p}检查非线性通量的增长、强制和单调性，追踪缩放后的方程与能量指数。
\item\label{b2:item:H2-goal-variation}在明确的边界类和源项对偶空间中构造变分解，并用比较、稳定性和径向点源检验结论的范围。
\end{itemize}
\end{learninggoals}

一块材料的导热系数可以跨界面跳跃，非线性介质的通量也可能随梯度大小变化。前一种问题不允许对系数随意微分，后一种问题则不允许把不同解相加后仍当作解。两者却都保留了弱形式：测试函数的梯度与通量配对，而通量沿解梯度的分量具有正的能量下界。以下从这个共同结构出发，分别说明还能取得哪些结论。

\ProseParagraphBreak
本附录只处理实值标量方程。\chapref{b2:ch:15}的能量解与\chapref{b2:ch:16}的可测系数局部估计提供先修；Sobolev 链式法则、Poincaré 不等式及临界嵌入按各处需要恢复具体形式。线性 Harnack 的证明还使用\secref{b2:sec:app-C-2}已证的局部 John--Nirenberg 指数尾。该结果由立方体停止分解建立，不依赖这里的椭圆正则性，所以这条证明路线没有循环。

\ProseParagraphBreak
两节的职责不同。第一节补全第十六章所述的线性正解比较；正幂迭代只能得到上界，必须再连接负幂平均才能控制下界。第二节从一般通量的局部能量估计转到一个有凸势的边值模型；局部上界只需增长和强制条件，唯一性另需比较两个梯度的单调性，极小化存在性又需要能量结构。\cite{Heinonen2006Nonlinear}提供非线性势论的进一步阅读方向，本附录不会把一般非线性 Harnack 或梯度正则性作为已经建立的结论。

% !TeX root = ../../Book2.tex
\section{能量估计与低正则迭代}
\label{b2:sec:H201}

系数只有可测性时，不能把方程逐项微分，却仍可以改变测试函数。截去低于某个高度的部分，可以观察高值集合；提高解的幂次，可以逐步增强可积性；取倒数或对数，则能接近正解的下界。这些操作都发生在弱形式的一阶导数之内。本节先把测试的合法性说清楚，再沿正幂、对数和倒数三条线建立线性 Harnack 比较。

\ProseParagraphBreak

以下讨论实标量方程。设 \(d\geq2\)，\(\Omega\subset\mathbb R^d\) 开，实矩阵场 \(A:\Omega\to\mathbb R^{d\times d}\) 可测，并满足
\[
 A(x)\xi\cdot\xi\geq\lambda|\xi|^2,
 \quad |A(x)\xi|\leq\Lambda|\xi|
 \quad(\xi\in\mathbb R^d)
\]
于几乎处处的 \(x\)，其中 \(0<\lambda\leq\Lambda<\infty\)。不要求 \(A\) 对称，也不使用系数的导数。实函数 \(u\in H^1_{\mathrm{loc}}(\Omega)\) 是齐次方程的弱次解，是指
\[
 \int_\Omega A\nabla u\cdot\nabla\varphi\leq0
 \quad(0\leq\varphi\in C_c^\infty(\Omega));
\]
若对全部实值 \(\varphi\in C_c^\infty(\Omega)\) 等式成立，则称为弱解。这个符号约定对应 \(-\operatorname{div}(A\nabla u)\leq0\)。下文的上、下界均先按几乎处处的意义理解，直到最后使用连续代表。

\subsection{Caccioppoli 不等式}
\label{b2:subsec:H20101}

把解本身或它的正部放到测试位置时，所得函数通常不光滑。弱形式可以接受这种测试，但理由来自 Sobolev 逼近以及通量的平方可积性，不能只靠形式上的分部积分。

\Needspace{6\baselineskip}
\begin{lemma}[非负 Sobolev 测试与正部]
\label{b2:lem:H2-sobolev-tests}
设 \(d\geq2\)，\(\Omega\subset\mathbb R^d\) 开，\(A:\Omega\to\mathbb R^{d\times d}\) 可测，且对几乎处处的 \(x\) 和全部 \(\xi\in\mathbb R^d\)，有
\[
 A(x)\xi\cdot\xi\geq\lambda|\xi|^2,
 \quad |A(x)\xi|\leq\Lambda|\xi|,
 \quad 0<\lambda\leq\Lambda<\infty.
\]
若实函数 \(u\in H^1_{\mathrm{loc}}(\Omega)\) 满足
\(\int_\Omega A\nabla u\cdot\nabla\varphi\leq0\) 对所有非负 \(\varphi\in C_c^\infty(\Omega)\) 成立，则同一不等式对每个非负、具有紧支集的 \(\varphi\in H^1(\Omega)\) 成立。对任意 \(k\in\mathbb R\)，函数 \((u-k)_+\) 也满足这一弱次解不等式。
\end{lemma}
\begin{proof}
先设非负 \(\varphi\in H^1(\Omega)\) 的支集包含于 \(K\Subset\Omega\)。在 \(K\) 附近取等于一的光滑紧支集函数，将 \(\varphi\) 乘以它再零延拓，便得到 \(H^1(\mathbb R^d)\) 函数。用非负磨光核卷积，得到 \(\varphi_\delta\geq0\)。当 \(\delta\) 足够小时，这些光滑函数的支集都落在同一个 \(U\Subset\Omega\) 内，并且 \(\varphi_\delta\to\varphi\) 于 \(H^1(U)\)。由于
\[
 \left|\int_U A\nabla u\cdot\nabla(\varphi_\delta-\varphi)\right|
 \leq\Lambda\|\nabla u\|_{L^2(U)}
       \cdot\|\nabla(\varphi_\delta-\varphi)\|_{L^2(U)},
\]
弱次解不等式可取极限。若原来是弱解，同一逼近不需要保留符号，因而允许任意实值的紧支撑 \(H^1\) 测试。

再令 \(v=(u-k)_+\)。Sobolev 截断链式法则给出
\(\nabla v=\mathbf1_{\{u>k\}}\cdot\nabla u\)。固定非负 \(\psi\in C_c^\infty(\Omega)\)，取分段线性函数
\[
 \theta_\delta(t)=
 \begin{cases}
 0,&t\leq0,\\
 t/\delta,&0<t<\delta,\\
 1,&t\geq\delta.
 \end{cases}
\]
函数 \(\theta_\delta(u-k)\cdot\psi\) 是合法的非负 Sobolev 测试。代入后得到
\[
 \int\theta_\delta(u-k)A\nabla u\cdot\nabla\psi
 +\int\psi\cdot\theta_\delta'(u-k)
                      A\nabla u\cdot\nabla u\leq0.
\]
第二项非负，故可从左边删除。第一项由紧支集上的
\(\Lambda|\nabla u|\cdot|\nabla\psi|\) 控制，这是可积函数；令 \(\delta\downarrow0\)，得到
\(\int A\nabla v\cdot\nabla\psi\leq0\)。分段函数折点上的导数取值不影响结果：Sobolev 函数在每个固定水平集上的梯度几乎处处为零，这也是截断链式法则的一部分。
\end{proof}

正部保留次解性质，使高度截断可以反复使用。下面的估计把弱形式中的椭圆性变成梯度控制；截断函数的导数记录了局部化所付的代价。

\Needspace{6\baselineskip}
\begin{proposition}[水平截断的 Caccioppoli 估计]
\label{b2:prop:H2-level-caccioppoli}
设 \(d\geq2\)，\(\Omega\subset\mathbb R^d\) 开，\(A:\Omega\to\mathbb R^{d\times d}\) 为可测实矩阵场，且
\(A(x)\xi\cdot\xi\geq\lambda|\xi|^2\)、\(|A(x)\xi|\leq\Lambda|\xi|\) 对几乎处处的 \(x\) 及全部 \(\xi\in\mathbb R^d\) 成立，其中 \(0<\lambda\leq\Lambda<\infty\)。设实函数 \(u\in H^1_{\mathrm{loc}}(\Omega)\) 满足
\(\int A\nabla u\cdot\nabla\varphi\leq0\) 对每个非负 \(\varphi\in C_c^\infty(\Omega)\) 成立。对 \(k\in\mathbb R\)、\(w=(u-k)_+\) 及 \(0\leq\eta\leq1\)、\(\eta\in C_c^\infty(\Omega)\)，有
\[
 \int_\Omega\eta^2|\nabla w|^2
 \leq4\frac{\Lambda^2}{\lambda^2}
          \int_\Omega w^2|\nabla\eta|^2
\]
以及
\begin{equation}\label{b2:eq:H2-level-energy}
 \int_\Omega|\nabla(\eta\cdot w)|^2
 \leq C\int_\Omega w^2|\nabla\eta|^2,
\end{equation}
其中 \(C\) 只依赖 \(\lambda,\Lambda\)。特别地，若
\(0<r<s\)、\(B_s=B(x_0,s)\Subset\Omega\)，则
\[
 \int_{B_r}|\nabla(u-k)_+|^2
 \leq\frac{C}{(s-r)^2}\int_{B_s}(u-k)_+^2.
\]
\end{proposition}
\begin{proof}
由上一引理，可用 \(\eta^2\cdot w\) 测试。在 \(w>0\) 处有 \(\nabla w=\nabla u\)，在其余部分 \(w=0\)、\(\nabla w=0\) 几乎处处，所以
\[
 \lambda\int\eta^2|\nabla w|^2
 \leq2\Lambda\int\eta w|\nabla w|\cdot|\nabla\eta|
 \leq\frac\lambda2\int\eta^2|\nabla w|^2
       +\frac{2\Lambda^2}{\lambda}\int w^2|\nabla\eta|^2.
\]
把梯度项移回左边即得第一式。再由
\[
 |\nabla(\eta\cdot w)|^2
 \leq2\eta^2|\nabla w|^2+2w^2|\nabla\eta|^2
\]
得到 \eqref{b2:eq:H2-level-energy}。最后取 \(\eta=1\) 于 \(B_r\)，支撑于 \(B_s\)，并使 \(|\nabla\eta|\leq C_d/(s-r)\)，即得两球估计。
\end{proof}

对齐次弱解，\(-u\) 也是解，因而负部可按相同方式估计；若改用 \(\eta^2\cdot(u-c)\) 测试，还可估计整个梯度。这一能量界不含 \(\nabla A\)，而\chapref{b2:ch:16}中的差商论证需要额外的系数正则性。两者控制的对象不同：这里从函数的局部大小控制一阶能量，差商方法则试图产生二阶弱导数。

\ProseParagraphBreak

齐次假设也有实际作用。若 \(-\operatorname{div}(A\nabla u)=f\)，测试会多出 \(\int f\eta^2w\)。必须先说明该积分有意义，再估计它；只有这样才能把源项的可积指数和缩放带入迭代。将这一项删除，会把含源项方程误当成齐次方程。

\subsection{De Giorgi 迭代}
\label{b2:subsec:H20102}

\cref{b2:thm:16-local-boundedness}已经对有界可测一致椭圆系数的实弱次解，完整证明了
\[
 \operatorname*{ess\,sup}_{B_R}u_+
 \leq C\left(R^{-d}\int_{B_{2R}}u_+^2\right)^{1/2}.
\]
该证明只使用水平截断能量、Sobolev 不等式及一个数列迭代，不需要 Harnack 比较。即使解只属于局部 \(H^1\)，也可先在稍小的外球上应用结论，再把外半径增大；若右侧有限，单调性给出所列形式。

\ProseParagraphBreak

这一过程称为 \term[DeGiorgiDiedai]{De Giorgi 迭代}[De Giorgi iteration]。它同时收缩半径、提高截断高度。若第 \(j\) 轮的高度为 \(k_j\)，则下一轮高值集合上有
\((u-k_j)_+\geq k_{j+1}-k_j\)，所以旧的能量还能控制新集合的测度。把这一测度界与 Sobolev 的积分增益结合，归一化后的能量满足
\[
 Y_{j+1}\leq Cb^jY_j^{1+\alpha},
 \quad b>1,\quad\alpha>0.
\]
起始高度由初始平方平均选择，使 \(Y_0\) 足够小，才会得到 \(Y_j\to0\)。超线性幂 \(1+\alpha\) 是小量改善的来源；若只有 \(Y_{j+1}\leq Cb^jY_j\)，这条论证不能迫使高值部分消失。

\ProseParagraphBreak

本节不重复这条已完成的高度迭代。下面改用幂函数，以便进一步接入倒数和对数测试。两种路线的区别也解释了方法的选择：障碍问题天然允许单侧高度截断；正解比较则需要把函数的大值与小值联系起来，幂函数的正、负方向更便于配合。

\subsection{Moser 迭代}
\label{b2:subsec:H20103}

逐次提高可积指数的路线称为 \term[MoserDiedai]{Moser 迭代}[Moser iteration]。形式上，用 \(u^{q-1}\) 测试会出现 \(u^{q-2}|\nabla u|^2\)，它正是 \(u^{q/2}\) 的梯度平方乘以一个常数。但本轮只知道 \(u\in L^q\)，并不知道这个幂已属于 \(H^1\)。因此必须先把幂函数截成全局 Lipschitz 函数，取得独立于截断高度的估计，再放开截断。

\ProseParagraphBreak

先记下本节使用的 Sobolev 指数：当 \(d>2\) 时取 \(\chi=d/(d-2)\)，当 \(d=2\) 时取 \(\chi=2\)。对 \(z\in H_0^1(B(0,2))\)，有
\(\|z\|_{2\chi}^2\leq C_d\|\nabla z\|_2^2\)。二维情形可先用零延拓后的 \(W^{1,4/3}\) 到 \(L^4\) 的 Sobolev 不等式，再用有限体积上的 Hölder 不等式把梯度的 \(L^{4/3}\) 范数界为 \(L^2\) 范数。缩放回任意 \(B_{2R}=B(x_0,2R)\)，所得形式是
\begin{equation}\label{b2:eq:H2-scaled-sobolev}
 \left(R^{-d}\int_{B_{2R}}|z|^{2\chi}\right)^{1/\chi}
 \leq C_dR^{2-d}\int_{B_{2R}}|\nabla z|^2
 \quad(z\in H_0^1(B_{2R})).
\end{equation}
这一个式子包含了两个维数范围。特别是在二维，若改写成未归一化的范数，右侧会带一个 \(R\)；不能把单位球常数原样用于任意半径。

\Needspace{6\baselineskip}
\begin{lemma}[幂函数能量估计]
\label{b2:lem:H2-power-energy}
设 \(d\geq2\)，\(x_0\in\mathbb R^d\)、\(R>0\)，记 \(B_t=B(x_0,t)\)。可测实矩阵场 \(A\) 在 \(B_{2R}\) 中满足
\(A(x)\xi\cdot\xi\geq\lambda|\xi|^2\)、\(|A(x)\xi|\leq\Lambda|\xi|\)，其中 \(0<\lambda\leq\Lambda<\infty\)，两式对几乎处处的 \(x\) 和全部 \(\xi\in\mathbb R^d\) 成立。设实函数 \(u\in H^1_{\mathrm{loc}}(B_{2R})\) 几乎处处非负，且
\(\int A\nabla u\cdot\nabla\varphi\leq0\) 对所有非负 \(\varphi\in C_c^\infty(B_{2R})\) 成立。取 \(q\geq2\)、\(R\leq r<s\leq2R\)，并假定 \(u\in L^q(B_s)\)。令 \(\chi=d/(d-2)\) 当 \(d>2\)，\(\chi=2\) 当 \(d=2\)。则
\[
 \left(R^{-d}\int_{B_r}u^{q\chi}\right)^{1/(q\chi)}
 \leq
 \left(\frac{Cq^2R^2}{(s-r)^2}\right)^{1/q}
 \left(R^{-d}\int_{B_s}u^q\right)^{1/q},
\]
其中 \(C\) 只依赖 \(d,\lambda,\Lambda\)。
\end{lemma}
\begin{proof}
固定 \(L>0\)。当 \(q>2\) 时，对 \(t\geq0\) 定义
\[
 h_L(t)=t\min(t,L)^{q-2},
 \quad w_L(t)=t\min(t,L)^{(q-2)/2};
\]
当 \(q=2\) 时直接定义 \(h_L(t)=w_L(t)=t\)。这样零点处没有 \(0^0\) 的约定问题。两函数在 \([0,\infty)\) 上 Lipschitz，零点取零，并可延拓为实轴上的 Lipschitz 函数。对于 \(q>2\)，其两支导数为
\[
 \begin{array}{c|cc}
 &0<t<L&t>L\\ \hline
 h_L'(t)&(q-1)t^{q-2}&L^{q-2}\\
 w_L'(t)&\dfrac q2t^{(q-2)/2}&L^{(q-2)/2}.
 \end{array}
\]
由这两行直接得到
\[
 h_L(t)^2\leq w_L(t)^2\cdot h_L'(t),
 \quad |w_L'(t)|^2\leq q^2\cdot h_L'(t),
 \quad h_L'(t)\geq0.
\]
第一式采用不含除法的形式，因而在 \(t=0\) 时仍有意义；\(q=2\) 时三式同样成立。折点 \(t=L\) 的导数值任取一侧即可，因为 \(\nabla u=0\) 几乎处处于 \(\{u=L\}\)。

取 \(0\leq\eta\leq1\)、\(\eta\in C_c^\infty(B_s)\)，且 \(\eta=1\) 于 \(B_r\)、\(|\nabla\eta|\leq C_d/(s-r)\)。对固定 \(L\)，函数 \(\eta^2\cdot h_L(u)\) 是合法的非负 Sobolev 测试。弱次解不等式给出
\[
 \begin{aligned}
 \lambda\int\eta^2h_L'(u)|\nabla u|^2
 &\leq2\Lambda\int\eta\cdot h_L(u)
                        |\nabla u|\cdot|\nabla\eta|\\
 &\leq\frac\lambda2\int\eta^2h_L'(u)|\nabla u|^2
       +\frac{2\Lambda^2}{\lambda}
                    \int w_L(u)^2|\nabla\eta|^2.
 \end{aligned}
\]
这里使用了 \(h_L\leq w_L\sqrt{h_L'}\)，再对两个非负因子使用 Young 不等式。吸收左边的一半，并用 \(|w_L'|^2\leq q^2h_L'\)，得到
\[
 \int|\nabla(\eta\cdot w_L(u))|^2
 \leq Cq^2\int w_L(u)^2|\nabla\eta|^2
 \leq\frac{Cq^2}{(s-r)^2}\int_{B_s}u^q.
\]
最后一步来自 \(w_L(u)^2\leq u^q\)。对 \(\eta\cdot w_L(u)\in H_0^1(B_{2R})\) 使用 \eqref{b2:eq:H2-scaled-sobolev}，得到
\[
 \left(R^{-d}\int_{B_r}w_L(u)^{2\chi}\right)^{1/\chi}
 \leq\frac{Cq^2R^2}{(s-r)^2}
              R^{-d}\int_{B_s}u^q.
\]
当 \(L\to\infty\) 时，\(w_L(u)\to u^{q/2}\) 几乎处处。左边用 Fatou 引理，右边保持不变，再取 \(q\) 次根即得结论。整个过程只使用本轮已有的 \(L^q\) 可积性，没有预设下一轮幂函数的梯度可积。
\end{proof}

\Needspace{6\baselineskip}
\begin{theorem}[正幂迭代的局部上界]
\label{b2:thm:H2-moser-upper}
设 \(d\geq2\)，\(x_0\in\mathbb R^d\)、\(R>0\)，记 \(B_t=B(x_0,t)\)。可测实矩阵场 \(A\) 在 \(B_{2R}\) 上满足
\(A(x)\xi\cdot\xi\geq\lambda|\xi|^2\)、\(|A(x)\xi|\leq\Lambda|\xi|\)，其中 \(0<\lambda\leq\Lambda<\infty\)，两式对几乎处处的 \(x\) 和全部实 \(\xi\) 成立。设实函数 \(u\in H^1_{\mathrm{loc}}(B_{2R})\cap L^2(B_{2R})\) 几乎处处非负，并满足
\(\int A\nabla u\cdot\nabla\varphi\leq0\) 对所有非负 \(\varphi\in C_c^\infty(B_{2R})\) 成立。则
\[
 \operatorname*{ess\,sup}_{B_R}u
 \leq C\left(R^{-d}\int_{B_{2R}}u^2\right)^{1/2},
\]
其中 \(C\) 只依赖 \(d,\lambda,\Lambda\)。
\end{theorem}
\begin{proof}
置 \(v(y)=u(x_0+Ry)\)、\(\widetilde A(y)=A(x_0+Ry)\)。在弱形式中换元，出现的共同正因子 \(R^{d-2}\) 可以约去；\(\widetilde A\) 保留同一组椭圆界。因此先证 \(R=1\)、\(x_0=0\) 的结论即可。

令 \(\chi\) 如幂函数能量估计所定，并取
\[
 q_j=2\chi^j,
 \quad r_j=1+2^{-j}
 \quad(j=0,1,2,\ldots).
\]
从 \(q_0=2\)、\(r_0=2\) 开始，上一引理逐次保证下一轮所需的可积性，而且
\[
 \|u\|_{L^{q_{j+1}}(B_{r_{j+1}})}
 \leq\left(Cq_j^2\cdot2^{2j+2}\right)^{1/q_j}
          \|u\|_{L^{q_j}(B_{r_j})}.
\]
把各轮相乘。因为
\[
 \sum_{j=0}^\infty
 \frac{\log C+2\log q_j+(2j+2)\log2}{q_j}<\infty,
\]
乘积被一个只依赖 \(d,\lambda,\Lambda\) 的常数控制。于是
\(\|u\|_{L^{q_j}(B_1)}\leq K\)，其中 \(K=C\|u\|_{L^2(B_2)}\)，与 \(j\) 无关。若存在 \(t>K\) 使 \(E=\{x\in B_1:u(x)>t\}\) 具有正测度，则
\[
 K\geq\|u\|_{L^{q_j}(B_1)}\geq t|E|^{1/q_j}
 \longrightarrow t,
\]
产生矛盾。因此 \(u\leq K\) 几乎处处。换元回原来的球，
\(\|v\|_{L^2(B_2)}=R^{-d/2}\|u\|_{L^2(B_{2R})}\)，即得所列估计。
\end{proof}

对有符号弱次解，可先由\cref{b2:lem:H2-sobolev-tests}取正部；对弱解再分别处理 \(u_+\) 和 \((-u)_+\)，便得到绝对值的局部上界。到这里，幂迭代已到达与 De Giorgi 迭代相同的结论。下一步需要把平方平均换成很小的正幂平均，因为对数的指数可积性只会产生一个小指数。

\Needspace{6\baselineskip}
\begin{lemma}[局部上界的低次平均版本]
\label{b2:lem:H2-small-power}
设 \(d\geq2\)，\(x_0\in\mathbb R^d\)、\(R>0\)，记 \(B_t=B(x_0,t)\)。可测实矩阵场 \(A\) 在 \(B_{2R}\) 中满足
\(A(x)\xi\cdot\xi\geq\lambda|\xi|^2\)、\(|A(x)\xi|\leq\Lambda|\xi|\)，其中 \(0<\lambda\leq\Lambda<\infty\)，两式对几乎处处的 \(x\) 和全部实 \(\xi\) 成立。实函数 \(u\in H^1_{\mathrm{loc}}(B_{2R})\) 几乎处处非负，并且
\(\int A\nabla u\cdot\nabla\varphi\leq0\) 对所有非负 \(\varphi\in C_c^\infty(B_{2R})\) 成立。对每个 \(a>0\)，有
\[
 \operatorname*{ess\,sup}_{B_R}u
 \leq C_a\left(R^{-d}\int_{B_{2R}}u^a\right)^{1/a},
\]
其中 \(C_a\) 只依赖 \(a,d,\lambda,\Lambda\)；右侧为无穷时按通常的扩展实数意义理解。
\end{lemma}
\begin{proof}
只需考虑右侧有限的情形，并缩放到 \(R=1\)。当 \(a\geq2\) 时，有限体积上的 Hölder 不等式保证 \(u\in L^2(B_2)\)，由上一项平方平均上界即得结论。

设 \(0<a<2\)，记 \(M(r)=\operatorname*{ess\,sup}_{B_r}u\)。首先，\(M(3/2)<\infty\)：用有限个较小球覆盖 \(\overline{B_{3/2}}\)，使每个球的二倍球都紧包含于 \(B_2\)，再逐球使用平方平均上界。所需平方积分来自 \(u\in H^1_{\mathrm{loc}}(B_2)\)。这个有限性先于后面的吸收步骤。

对 \(1\leq r<s\leq3/2\)，在 \(\overline{B_r}\) 各点放置半径 \(\rho=(s-r)/4\) 的球，其二倍球均在 \(B_s\) 内。取有限子覆盖，在每个球上使用已证估计，得
\[
 \begin{aligned}
 M(r)&\leq C(s-r)^{-d/2}\|u\|_{L^2(B_s)}\\
 &\leq C(s-r)^{-d/2}M(s)^{1-a/2}
               \left(\int_{B_{3/2}}u^a\right)^{1/2}.
 \end{aligned}
\]
若 \(M(s)=0\)，结论已成立。否则令
\(U=(\int_{B_{3/2}}u^a)^{1/a}\)，对上式用指数 \(2/(2-a)\) 与 \(2/a\) 的带参数 Young 不等式。即使 \(a<1\)，这个积分幂量仍有意义，且此处不使用范数三角不等式。对任意 \(0<\vartheta<1\)，得到
\[
 M(r)\leq\vartheta M(s)
       +C_{a,\vartheta}(s-r)^{-d/a}U.
\]
取 \(r_j=3/2-2^{-j-1}\)，于是 \(r_0=1\)，且
\(r_{j+1}-r_j=2^{-j-2}\)。反复代入可得
\[
 M(1)\leq\vartheta^nM(r_n)
   +C_{a,\vartheta}2^{2d/a}U
        \sum_{j=0}^{n-1}\left(\vartheta2^{d/a}\right)^j.
\]
选择 \(\vartheta<2^{-d/a}\)，级数可和；由于
\(M(r_n)\leq M(3/2)<\infty\)，第一项趋零。故
\(M(1)\leq C_aU\)，扩大右侧积分域到 \(B_2\)，再恢复尺度即得结论。这里没有把未知的无穷上确界移到不等式另一边。
\end{proof}

仅有正幂估计仍然不能给出正下界。对正解取对数后，加法常数对应原函数的乘法大小，而平均振荡消去了这个常数。因此，对数能量适合连接正幂与负幂：先控制与整体高度无关的振荡，再用指数可积性把它还原成两个方向的平均。

\Needspace{6\baselineskip}
\begin{proposition}[对数能量与正负平均连接]
\label{b2:prop:H2-log-bridge}
设 \(d\geq2\)，\(\Omega\subset\mathbb R^d\) 开，\(A:\Omega\to\mathbb R^{d\times d}\) 为可测实矩阵场，且
\(A(x)\xi\cdot\xi\geq\lambda|\xi|^2\)、\(|A(x)\xi|\leq\Lambda|\xi|\) 对几乎处处的 \(x\) 和全部实 \(\xi\) 成立，其中 \(0<\lambda\leq\Lambda<\infty\)。设实函数 \(u\in H^1_{\mathrm{loc}}(\Omega)\) 几乎处处非负，且
\(\int A\nabla u\cdot\nabla\varphi=0\) 对每个实值 \(\varphi\in C_c^\infty(\Omega)\) 成立。若 \(Q\) 为轴平行立方体，\(4Q\Subset\Omega\)，则存在只依赖 \(d,\lambda,\Lambda\) 的 \(a_0>0\)，使每个 \(\varepsilon>0\)、\(v=u+\varepsilon\) 满足
\[
 \left(\frac1{|Q|}\int_Qv^{a_0}\right)
 \left(\frac1{|Q|}\int_Qv^{-a_0}\right)\leq9.
\]
此外，对每个实值 \(\eta\in C_c^\infty(\Omega)\)，有
\[
 \int_\Omega\eta^2|\nabla\log v|^2
 \leq4\frac{\Lambda^2}{\lambda^2}
                  \int_\Omega|\nabla\eta|^2.
\]
这里 \(4Q\) 表示与 \(Q\) 同中心、边长为四倍的立方体。
\end{proposition}
\begin{proof}
常数平移不改变梯度，所以 \(v\) 仍为齐次弱解。因为 \(v\geq\varepsilon\)，倒数函数在其取值范围内及其一阶导数都有界，\(\eta^2/v\) 属于紧支撑的 \(H^1\)，可作为测试。又有 \(\log v\in H^1_{\mathrm{loc}}(\Omega)\)：负部由 \(\varepsilon\) 控制，正部可由 \(u\) 加一个依赖 \(\varepsilon\) 的常数控制，而链式法则给出 \(\nabla\log v=\nabla u/v\)。弱方程因此给出
\[
 \begin{aligned}
 \lambda\int\eta^2|\nabla\log v|^2
 &\leq2\Lambda\int|\eta|\cdot|\nabla\log v|
                                  \cdot|\nabla\eta|\\
 &\leq\frac\lambda2\int\eta^2|\nabla\log v|^2
                  +\frac{2\Lambda^2}{\lambda}
                             \int|\nabla\eta|^2.
 \end{aligned}
\]
吸收后得到对数能量界，常数与 \(\varepsilon\) 无关。

令 \(b=\log v\)。对 \(Q\) 的任意轴平行子立方体 \(P\)，包括 \(P=Q\)，取在 \(P\) 上为一、支撑于 \(2P\) 的截断，且梯度不超过 \(C_d/\ell(P)\)。所有 \(2P\) 都包含于 \(3Q\)，所以这些测试统一合法。能量界给出
\[
 \int_P|\nabla b|^2\leq C\ell(P)^{d-2}.
\]
立方体的 Poincaré 不等式和 Cauchy--Schwarz 不等式于是给出
\[
 \begin{aligned}
 \frac1{|P|}\int_P|b-b_P|
 &\leq\left(\frac1{|P|}\int_P|b-b_P|^2\right)^{1/2}\\
 &\leq C_d\ell(P)
          \left(\frac1{|P|}\int_P|\nabla b|^2\right)^{1/2}
 \leq M,
 \end{aligned}
\]
其中可选 \(M>0\) 只依赖 \(d,\lambda,\Lambda\)。注意，所得界控制了 \(Q\) 内每个子立方体，而不只是 \(Q\) 本身；这正是下一步所需的信息。

现在使用\cref{b2:cor:C-local-john-nirenberg}。其假设是初始立方体内的局部可积性，以及所有二进子立方体上的同一平均振荡界。这里 \(b\in H^1(Q)\subset L^1(Q)\)，而刚证得的界对全部轴平行子立方体成立，因而满足该假设。取推论中的维数常数 \(c_d>0\)，令
\(a_0=c_d/(2M)\)，便有
\[
 \frac1{|Q|}\int_Qe^{\pm a_0(b-b_Q)}\leq3.
\]
分别乘回 \(e^{a_0b_Q}\) 和 \(e^{-a_0b_Q}\)，再将两个平均相乘，未知常数因子抵消，得到所列乘积界。附录 C 的局部推论从停止立方体和指数尾的积分直接得到；它的证明不使用 Harnack 比较，因此此处没有循环依赖。
\end{proof}

有了正负平均的连接，只差把负幂平均变成下界。最方便的做法是对倒数使用已证的次解上界，而不是重新进行一遍负幂迭代。

\Needspace{6\baselineskip}
\begin{theorem}[线性正解比较的补充完整证明]
\label{b2:thm:H2-linear-harnack}
设 \(d\geq2\)，\(x_0\in\mathbb R^d\)、\(R>0\)，记 \(B_t=B(x_0,t)\)。可测实矩阵场 \(A\) 在 \(B_{2R}\) 中满足
\(A(x)\xi\cdot\xi\geq\lambda|\xi|^2\)、\(|A(x)\xi|\leq\Lambda|\xi|\)，其中 \(0<\lambda\leq\Lambda<\infty\)，两式对几乎处处的 \(x\) 和全部 \(\xi\in\mathbb R^d\) 成立。若实函数 \(u\in H^1_{\mathrm{loc}}(B_{2R})\) 几乎处处非负，并且
\(\int A\nabla u\cdot\nabla\varphi=0\) 对所有实值 \(\varphi\in C_c^\infty(B_{2R})\) 成立，则
\[
 \operatorname*{ess\,sup}_{B_R}u
 \leq C\operatorname*{ess\,inf}_{B_R}u,
\]
其中 \(C\) 只依赖 \(d,\lambda,\Lambda\)。
\end{theorem}
\begin{proof}
先固定 \(\varepsilon>0\)，令 \(v=u+\varepsilon\)。其倒数 \(w=v^{-1}\) 属于局部 \(H^1\)，而且是非负弱次解。为确认不等式方向，在 \(v\) 的弱方程中使用
\(v^{-2}\cdot\varphi\)，其中 \(0\leq\varphi\in C_c^\infty(B_{2R})\)。这个测试合法，因为 \(v\geq\varepsilon\)。展开导数，得到
\[
 \begin{aligned}
 \int A\nabla w\cdot\nabla\varphi
 &=-\int v^{-2}A\nabla v\cdot\nabla\varphi\\
 &=-2\int v^{-3}\varphi A\nabla v\cdot\nabla v\leq0.
 \end{aligned}
\]
最后一项可积，因为 \(v^{-3}\) 有界、\(\nabla v\) 局部平方可积；其符号只使用椭圆性，不要求矩阵对称。

取 \(4Q\Subset B_{2R}\)。低次平均上界可从球转到立方体：在 \(\overline{Q/2}\) 各点放置半径 \(\ell(Q)/16\) 的球，二倍球仍在 \(Q\) 内，并选有限子覆盖。逐球使用\cref{b2:lem:H2-small-power}，球尺度因子与 \(|Q|^{-1}\) 只相差一个维数常数。因此，对上一个命题给出的同一个 \(a_0>0\)，有
\[
 \begin{aligned}
 \operatorname*{ess\,sup}_{Q/2}v
 &\leq C\left(\frac1{|Q|}\int_Qv^{a_0}\right)^{1/a_0},\\
 \operatorname*{ess\,sup}_{Q/2}v^{-1}
 &\leq C\left(\frac1{|Q|}\int_Qv^{-a_0}\right)^{1/a_0}.
 \end{aligned}
\]
这里 \(v\) 和 \(v^{-1}\) 都是非负次解。相乘后使用对数连接，得到
\[
 \operatorname*{ess\,sup}_{Q/2}v
 \leq H\operatorname*{ess\,inf}_{Q/2}v,
\]
其中 \(H\geq1\) 只依赖 \(d,\lambda,\Lambda\)。左边乘积可以改写成上、下界之比，是因为 \(v\geq\varepsilon\)，且它在 \(Q/2\) 上已有有限上界。

下面把这些小立方体的比较接成整球比较。缩放到 \(R=1\)、\(x_0=0\)，取
\(\ell=1/(10\sqrt d)\)、\(h=\ell/16\)。在 \(\overline{B_1}\) 选一个极大的有限 \(h\)-分离点集 \(c_1,\ldots,c_N\)。极大性说明每个点到某个 \(c_i\) 的距离不超过 \(h\)；而两两不交的球 \(B(c_i,h/2)\) 都在 \(B_{1+h/2}\) 内，所以
\(N\leq(1+2/h)^d\)。令 \(Q_i\) 为中心 \(c_i\)、边长 \(\ell\) 的轴平行立方体。各 \(4Q_i\) 都紧包含于 \(B_2\)，各 \(E_i=Q_i/2\) 覆盖 \(\overline{B_1}\)。

把距离不超过 \(3h\) 的两个中心连成边。所得有限图连通：任取两个中心，沿它们之间位于 \(\overline{B_1}\) 的线段，以不超过 \(h\) 的步长取点，再把每个取样点换成距离不超过 \(h\) 的某个中心；相邻中心的距离便不超过 \(3h\)。每一条边对应的 \(E_i,E_j\) 有正测度重叠，因为
\(3h<\ell/2\)。记
\(S_i=\operatorname*{ess\,sup}_{E_i}v\)、
\(m_i=\operatorname*{ess\,inf}_{E_i}v\)。各片上有 \(S_i\leq Hm_i\)，相邻片的重叠又给出 \(m_i\leq S_j\)。于是沿一条不重复顶点的链，得到
\(S_i\leq H^Nm_j\)。这个论证只比较几乎处处的上下界；有限多个零测例外的并仍为零测集。

选择使 \(S_i\) 最大和使 \(m_j\) 最小的两片，再利用覆盖，便得
\[
 \operatorname*{ess\,sup}_{B_1}v
 \leq\max_iS_i
 \leq H^N\min_jm_j
 \leq H^N\operatorname*{ess\,inf}_{B_1}v.
\]
因为 \(N\) 只受维数控制，常数仍只依赖 \(d,\lambda,\Lambda\)。恢复尺度后，令 \(\varepsilon\downarrow0\)。本性上确界和本性下确界在加常数时都平移同一常数，因此得到 \(u\) 的结论，即使其下确界为零也不例外。
\end{proof}

证明中的加常数步骤避免了预先假设解严格为正。它还显示，正幂上界、倒数次解和对数振荡承担不同任务：前者给局部最大值，倒数把同一个结论转成最小值，对数连接则保证两边平均可以相消。只完成其中的正幂迭代，还不足以得到 Harnack 不等式。

\subsection{De Giorgi--Nash--Moser 理论与 Harnack 不等式}
\label{b2:subsec:H20104}

本节补全了\cref{b2:thm:16-measurable-harnack}中的线性比较路线。相关公开论文 \cite[Theorems 1.1--1.3 and Section 2]{Li2017Harnack}讨论了局部上界、正性传播与 Harnack 比较之间的关系；这里实际使用的指数尾则由\appref{b2:app:C}在书内证明。上述论证始终保留了实标量、非负解、内球和一致椭圆假设。

\ProseParagraphBreak

Harnack 比较对有符号解的作用，要通过减去常数来显现。在外球上记 \(M=\operatorname*{ess\,sup}u\)、\(m=\operatorname*{ess\,inf}u\)，则 \(u-m\) 和 \(M-u\) 是两个非负弱解。分别应用比较并相加，便得到较小球上的严格振荡收缩；\subsecref{b2:subsec:160404}已经给出收缩因子、二进半径迭代以及从本性上下界构造 Hölder 连续代表的完整步骤。现在该推导所需的比较也已在本节证明，无须再把连续性当作额外输入。

\ProseParagraphBreak

正解比较还有一个可以直接读出的后果：非负解若在某个小球上具有零的本性下界，就必须在该球上消失；连通性将这种消失传到整个区域。这是标量方程中强最大值原理的一种内部形式。

\Needspace{6\baselineskip}
\begin{corollary}[连通区域中的正性与强最大值原理]
\label{b2:cor:H2-connected-positivity}
设 \(d\geq2\)，\(\Omega\subset\mathbb R^d\) 为连通开集，\(A:\Omega\to\mathbb R^{d\times d}\) 为可测实矩阵场，并满足
\(A(x)\xi\cdot\xi\geq\lambda|\xi|^2\)、\(|A(x)\xi|\leq\Lambda|\xi|\) 对几乎处处的 \(x\) 和全部实 \(\xi\) 成立，其中 \(0<\lambda\leq\Lambda<\infty\)。若实函数 \(u\in H^1_{\mathrm{loc}}(\Omega)\) 几乎处处非负，且
\(\int A\nabla u\cdot\nabla\varphi=0\) 对全部实值 \(\varphi\in C_c^\infty(\Omega)\) 成立，则或者 \(u=0\) 几乎处处，或者对每个 \(B_{2r}(x)\Subset\Omega\)，都有
\[
 \operatorname*{ess\,inf}_{B_r(x)}u>0.
\]
因而，非零解的局部 Hölder 连续代表在 \(\Omega\) 中处处为正。更一般地，若实弱解 \(z\in H^1_{\mathrm{loc}}(\Omega)\) 满足 \(z\leq M\) 几乎处处，其中 \(M\in\mathbb R\)，且其连续代表在某个内点取值 \(M\)，则 \(z=M\) 几乎处处，其连续代表处处等于 \(M\)。
\end{corollary}
\begin{proof}
令 \(Z\) 为所有具有一个邻域、使 \(u\) 在该邻域上几乎处处为零的点。按定义，\(Z\) 是相对开集。它也相对闭：若 \(x\in\overline Z\cap\Omega\)，取 \(B_{2r}(x)\Subset\Omega\)。球 \(B_r(x)\) 与 \(Z\) 的交集含有非空开集，所以 \(u\) 在一个正测度子集上为零；由非负性，\(\operatorname*{ess\,inf}_{B_r(x)}u=0\)。Harnack 比较给出 \(u=0\) 几乎处处于 \(B_r(x)\)，故 \(x\in Z\)。连通性说明 \(Z\) 为空或等于 \(\Omega\)；后一情形用可数子覆盖合并零测例外，得到 \(u=0\) 几乎处处于整个区域。

若 \(u\) 不是零函数，\(Z\) 必为空。任意 \(B_{2r}(x)\Subset\Omega\) 的内球若有零的本性下界，Harnack 比较会使该内球成为 \(Z\) 的一部分，矛盾。因此本性下界严格为正。连续代表与原函数几乎处处相等，故在内球上也不能低于这一正下界；否则连续性会给出一个正测度的违反集合。

最后对 \(u=M-z\) 应用这个二择结论。它是非负弱解，其连续代表在指定内点为零，因而不可能属于处处严格为正的情形。于是 \(z=M\) 几乎处处；连续性再把等式延伸到每一点。
\end{proof}

上述局部有界、正解比较和 Hölder 正则性通常合称 \term[DeGiorgiNashMoserLilun]{De Giorgi--Nash--Moser 理论}[De Giorgi--Nash--Moser theory]。名字相连并不意味着三个结论可以互相省略。对于椭圆系统、复值方程或一般非齐次方程，标量正部测试、序结构或齐次缩放可能失效。下一节转向 \(p\) 次增长方程时，会重新建立相应的截断估计和局部上界，而不把本节的线性 Harnack 定理直接套入。

% !TeX root = ../../Book2.tex
\section{非线性通量与变分解}
\label{b2:sec:H202}

当通量是梯度的非线性函数时，线性叠加一般不再成立，凸性和单调性取代了许多线性代数工具。本节固定 \(1<p<\infty\)，从一般 \(p\) 次增长的局部估计开始，再对最基本的凸能量完整证明边值解的存在、唯一与比较。

\subsection{拟线性椭圆方程}
\label{b2:subsec:H20201}

设 \(a(x,\xi)\) 对 \(x\) 可测、对 \(\xi\) 连续，且几乎处处满足
\[a(x,\xi)\cdot\xi\ge\lambda|\xi|^p,
 \qquad |a(x,\xi)|\le\Lambda|\xi|^{p-1}.\]
这类 \term[NixianxingTuoyuanFangcheng]{拟线性椭圆方程}[quasilinear elliptic equation] 的散度型弱形式为
\[\int_\Omega a(x,\nabla u)\cdot\nabla\varphi=0,
 \qquad\varphi\in C_c^\infty(\Omega),\quad u\in W^{1,p}_{\mathrm{loc}}(\Omega).\]
增长条件使通量局部属于 \(L^{p'}\)，其中 \(p'=p/(p-1)\)，所以积分有意义并可延拓到紧支撑的 \(W^{1,p}\) 测试。这里的条件不含 \(u\) 自身或低阶项；更一般方程须为新增项另列增长与可积性假设。

\begin{proposition}[非线性截断估计]
\label{b2:prop:H2-p-caccioppoli}
若 \(u\) 为上述方程的弱次解，\(w=(u-k)_+\)，则任意 \(0\le\eta\le1\) 光滑紧支集截断满足
\[\int\eta^p|\nabla w|^p\le C\int w^p|\nabla\eta|^p,
 \qquad \int|\nabla(\eta w)|^p\le C\int w^p|\nabla\eta|^p.\]
\end{proposition}
\begin{proof}
用 \(\eta^pw\) 测试。梯度仅在 \(u>k\) 出现，主项由强制条件控制为 \(\lambda\int\eta^p|\nabla w|^p\)，交叉项的绝对值不超过
\[p\Lambda\int(\eta|\nabla w|)^{p-1}w|\nabla\eta|.\]
对共轭指数 \(p',p\) 使用带参数的 Young 不等式，吸收主项的一半，得到第一式；再由 \(|\nabla(\eta w)|^p\le C_p(\eta^p|\nabla w|^p+w^p|\nabla\eta|^p)\) 得第二式。若 \(k<0\)，论证仍成立，因为方程通量不依赖 \(u\) 的值。
\end{proof}

\begin{theorem}[非线性次解的局部有界性]
\label{b2:thm:H2-p-bounded}
上述次解在 \(B_{2R}\) 中满足
\[\operatorname*{ess\,sup}_{B_R}u_+\le C\left(R^{-d}\int_{B_{2R}}u_+^p\right)^{1/p}.\]
常数只依赖 \(d,p,\lambda,\Lambda\)。
\end{theorem}
\begin{proof}
先缩放至单位球；变换后的通量保持相同增长常数。若 \(p<d\)，取 \(\chi=d/(d-p)>1\)，若 \(p\ge d\)，取任意固定 \(\chi>1\)，例如 \(2\)。对紧支撑于 \(B_2\) 的 \(z\)，均有 \(\lVert z\rVert_{p\chi}\le C\lVert\nabla z\rVert_p\)：临界及超临界情形可选择 \(q<d\) 使 \(q^*=p\chi\)，再用有界球上 \(L^p\) 控制 \(L^q\)。一维直接由积分基本定理控制上确界，以下只需考虑 \(d\ge2\)。

设 \(r_j=1+2^{-j}\)、\(k_j=K(1-2^{-j})\)、\(Y_j=\int_{B_{r_j}}(u-k_j)_+^p\)。在两相邻球间取截断，并应用上个命题和 Sobolev 不等式。由于集合 \(E_{j+1}=B_{r_{j+1}}\cap\{u>k_{j+1}\}\) 的测度至多为 \((k_{j+1}-k_j)^{-p}Y_j\)，Hölder 不等式给出
\[Y_{j+1}\le C2^{pj}Y_j\cdot|E_{j+1}|^{1-1/\chi}
 \le Cb^jK^{-p\alpha}Y_j^{1+\alpha},\qquad\alpha=1-1/\chi>0,\]
其中 \(b>1\) 只依赖固定指数。令 \(Z_j=Y_j/K^p\)，则 \(Z_{j+1}\le Cb^jZ_j^{1+\alpha}\)。若 \(Z_0\le C^{-1/\alpha}b^{-1/\alpha^2}\)，归纳可得 \(Z_j\le C^{-1/\alpha}b^{-j/\alpha-1/\alpha^2}\)，从而趋零。选 \(K\) 为 \(\lVert u_+\rVert_{L^p(B_2)}\) 的固定倍数即可满足该条件。极限说明 \((u-K)_+=0\) 于 \(B_1\)，最后恢复空间尺度即得结论。
\end{proof}

这里没有使用通量的严格单调性；局部能量与增长已经足以建立上界。唯一性却需要比较两个不同梯度，只有单点强制条件无法保证，因此须增加下面的单调结构。

\subsection{\texorpdfstring{\(p\)}{p}-Laplace 方程}
\label{b2:subsec:H20202}

基本模型为 \term[pLaplaceSuanzi]{\(p\)-Laplace 算子}[p-Laplace operator]
\[\Delta_pu\coloneqq\operatorname{div}(|\nabla u|^{p-2}\nabla u).\]
在零梯度处把通量定义为零；\(p>1\) 保证这个延拓连续。记 \(F(\xi)=|\xi|^{p-2}\xi\)，它是严格凸函数 \(|\xi|^p/p\) 的梯度。

\begin{lemma}[通量的严格单调性]
\label{b2:lem:H2-monotonicity}
若 \(\xi\ne\eta\)，则 \((F(\xi)-F(\eta))\cdot(\xi-\eta)>0\)。更具体地，存在 \(c_p>0\) 使
\[(F(\xi)-F(\eta))\cdot(\xi-\eta)\ge
 \begin{cases}c_p|\xi-\eta|^p,&p\ge2,\\
 c_p(|\xi|+|\eta|)^{p-2}|\xi-\eta|^2,&1<p<2.
 \end{cases}\]
第二式在两向量均为零时理解为零。
\end{lemma}
\begin{proof}
令 \(h=\xi-\eta\)。对非零 \(z\)，矩阵 \(DF(z)\) 在径向及切向的特征值分别为 \((p-1)|z|^{p-2}\) 与 \(|z|^{p-2}\)。沿线段积分得到
\[(F(\xi)-F(\eta))\cdot h\ge\min(1,p-1)|h|^2\int_0^1|\eta+th|^{p-2}\,\mathrm dt.\]
线段穿过零点时可能出现 \(|t-t_0|^{p-2}\)，它因 \(p>1\) 可积，故可分段取极限；\(h=0\) 情形显然。若 \(p<2\)，\(|\eta+th|\le|\xi|+|\eta|\) 且指数为负，直接得到所列下界。若 \(p\ge2\)，投影到 \(h\) 的方向，\(|\eta+th|\ge|h||t+c|\)。区间长度为一时，\(\int_0^1|t+c|^{p-2}\,\mathrm dt\) 的最小值在中心对称位置达到，等于 \(2^{2-p}/(p-1)\)。代入即得正的常数倍 \(|h|^p\)。
\end{proof}

当 \(p=2\) 时通量线性，回到 Laplace 方程。\(p>2\) 时零梯度处线性化退化，\(1<p<2\) 时其系数在零梯度附近奇异。因此，不能先形式微分，再把所得方程当成具有统一有界系数的线性方程。即便解已经局部有界，梯度的上、下界仍未由此获得。

\subsection{非线性变分问题}
\label{b2:subsec:H20203}

设 \(\Omega\) 有界，\(a\) 可测且 \(0<a_0\le a(x)\le a_1\)。令 \(g\in W^{1,p}(\Omega)\)，\(X=W_0^{1,p}(\Omega)\) 以梯度范数赋范，\(\ell\in X^*\)。在仿射类 \(g+X\) 上定义
\[J(g+w)\coloneqq\frac1p\int_\Omega a(x)|\nabla(g+w)|^p\,\mathrm dx-\ell(w).\]
这里仅对零边界增量施加泛函，因此不必先把 \(\ell\) 延拓到整个 \(W^{1,p}\)。不同延拓只改变固定边界类上的常数，不改变极小元。

\begin{theorem}[加权 \(p\) 次能量的 Dirichlet 解]
\label{b2:thm:H2-variational}
上述能量有唯一极小元 \(u\in g+X\)，它是该边界类中唯一满足
\[\int_\Omega a(x)|\nabla u|^{p-2}\nabla u\cdot\nabla\varphi=\ell(\varphi)
 \quad(\varphi\in X)\]
的弱解。若 \(g=0\)，还有 \(\lVert\nabla u\rVert_p^{p-1}\le a_0^{-1}\lVert\ell\rVert_{X^*}\)。
\end{theorem}
\begin{proof}
由 \(|\nabla w|^p\le2^{p-1}(|\nabla(g+w)|^p+|\nabla g|^p)\)，能量下界为 \(c\lVert\nabla w\rVert_p^p-\lVert\ell\rVert\lVert\nabla w\rVert_p-C_g\)，故极小序列在 \(X\) 有界。\(1<p<\infty\) 保证自反性，因而可取弱收敛子列。积分项为凸泛函，且强连续：基本不等式
\[\bigl||\xi|^p-|\eta|^p\bigr|\le p(|\xi|+|\eta|)^{p-1}|\xi-\eta|\]
与 Hölder 不等式给出该连续性。凸强连续泛函弱下半连续，连续线性项弱连续，所以第十四章直接法的论证给出极小元。

对任意 \(\varphi\in X\)，能量沿 \(u+t\varphi\) 可微。差商由常数倍 \((|\nabla u|+|\nabla\varphi|)^{p-1}|\nabla\varphi|\) 控制，此函数可积；控制收敛给出所列弱等式。反之，凸函数的切平面不等式说明满足弱等式者使能量全局最小。若有两个解，以差测试两式之差，上一引理及 \(a\ge a_0\) 迫使梯度相等，零迹差的 Poincaré 不等式给出函数相等。最后在零边界情形以 \(u\) 测试，有 \(a_0\lVert\nabla u\rVert_p^p\le\lVert\ell\rVert\lVert\nabla u\rVert_p\)，约去即可。
\end{proof}

同一单调性也给出比较原理。若两解边界有序、源项有序，用正部差测试即可推出解有序，完整计算见本附录习题。存在性证明使用了能量的势结构；一般单调通量未必等于某个势的梯度，不能因此宣称上面的极小化证明已经覆盖全部拟线性方程。

\subsection{径向解与低正则理论的边界}
\label{b2:subsec:H20204}

设 \(u(x)=U(r)\)、\(r=|x|>0\)。径向散度计算给出
\[\Delta_pu=r^{1-d}\frac{\mathrm d}{\mathrm dr}\bigl(r^{d-1}|U'|^{p-2}U'\bigr).\]
于是穿孔区域的非恒定径向 \(p\)-调和函数为
\[U(r)=A+Br^{(p-d)/(p-1)}\quad(p\ne d),\qquad
 U(r)=A+B\log r\quad(p=d).\]
计算从恒定径向通量开始，先解 \(|U'|^{p-2}U'=Cr^{1-d}\)，再积分即可。公式中的常数 \(B\) 已吸收指数与通量之间的非线性关系，不能把 \(B\) 直接认作源点质量。

\ProseParagraphBreak

这个例子还区分了“穿孔域中解方程”和“在原点两侧弱解方程”。非零径向通量穿过每个小球面的总量为同一非零常数；即使 \(p>d\) 时函数在原点有有限极限，它在全球中仍携带点源。此时梯度的 \(p\) 次径向积分为常数倍 \(\int_0^1r^{-(d-1)/(p-1)}\,\mathrm dr\)，恰因 \(p>d\) 而有限，所以这是具有点源的 Sobolev 解，不能把点源忽略。\(p\le d\) 时该能量发散，与同一函数不属于全球自然能量空间一致。

\ProseParagraphBreak

对于一般 \(p\) 次增长方程，进一步的正解 Harnack 不等式、局部 Hölder 连续及更精细的梯度正则性需要额外论证。这里已经证明截断能量、局部有界和凸势模型的边值存在唯一性，尚未用线性 Harnack 代替这些非线性步骤。阅读非线性势论时，应逐项查看方程是否齐次、是否有低阶项、右端是函数还是测度，以及结论针对解本身还是梯度；这些差别会改变自然尺度与有效指数。


\begin{chaptersummary}
弱形式中的截断测试先把通量的正性变成局部能量界，再由 Sobolev 增益推动迭代。在线性齐次方程中，正幂上界、对数能量、指数尾和倒数次解各有职责：它们依次提供上界、平均振荡控制、正负平均连接和下界，最终得到不要求系数可微的 Harnack 比较。

\ProseParagraphBreak
非线性问题还须分清局部估计、边值可解性与精细正则性。\(p\) 次增长支持截断迭代，凸能量支持直接法，定量单调性则控制唯一性与源项稳定性。径向通量和点值测试说明，原点处连续不等于原点处齐次：点源的强度由球面通量决定，而它能否进入自然对偶空间由 \(p\) 与维数的关系决定。一般非线性正解比较和梯度理论仍需在这些工具之后继续建立。
\end{chaptersummary}

\BookExercises

以下八题为本书自拟，分别考察边界类、缩放、障碍约束、点源以及非线性解算子的性质。题中的函数、测试函数和对偶泛函均取实值。

% 自拟：径向通量的边界归一化，并以迹和严格单调性排除非径向解。
\Needspace{6\baselineskip}
\begin{exercise}\label{b2:ex:H2-annulus}
设 \(d\ge2\)、\(1<p<\infty\)、\(0<a<b\)，并令
\(\Omega=\{\,x\in\mathbb R^d\mid a<|x|<b\,\}\)。给定 \(A,B\in\mathbb R\)，求 \(\Omega\) 中径向的弱 \(p\)-调和函数，其内、外边界迹分别为常数 \(A,B\)。证明所得函数在所有具有这两个边界迹的 \(W^{1,p}(\Omega)\) 弱解中唯一；竞争者不必径向。
\end{exercise}
\begin{solution}
径向解 \(u(x)=U(|x|)\) 的通量守恒要求
\(r^{d-1}|U'(r)|^{p-2}U'(r)\) 为常数。因此先令
\[
 I=\int_a^b t^{-(d-1)/(p-1)}\,\mathrm dt,
 \quad
 U(r)=A+\frac{B-A}{I}\int_a^r t^{-(d-1)/(p-1)}\,\mathrm dt.
\]
这里 \(I>0\)，且 \(U(a)=A\)、\(U(b)=B\)。当 \(p\ne d\) 时，记 \(\gamma=(p-d)/(p-1)\)，积分给出
\[
 U(r)=A+(B-A)\frac{r^\gamma-a^\gamma}{b^\gamma-a^\gamma}.
\]
当 \(p=d\) 时，相应公式为
\[
 U(r)=A+(B-A)\frac{\log(r/a)}{\log(b/a)}.
\]
这些公式也包含 \(A=B\) 的常值情形。

令 \(c=(B-A)/I\)。由 \(U'(r)=cr^{-(d-1)/(p-1)}\)，有
\[
 r^{d-1}|U'(r)|^{p-2}U'(r)=|c|^{p-2}c.
\]
环域离开原点，故 \(u\) 在闭环域上光滑，属于 \(W^{1,p}(\Omega)\)，其迹就是上述边界值。非恒定情形的梯度不为零，径向散度公式给出 \(\Delta_pu=0\)；常值情形的通量恒为零。对紧支撑光滑测试分部积分，再利用通量属于 \(L^{p'}(\Omega)\) 和测试函数的稠密性，即得对全部 \(\varphi\in W_0^{1,p}(\Omega)\) 的弱等式。

若 \(v\in W^{1,p}(\Omega)\) 是另一弱解且迹相同，光滑环域上的零迹刻画保证 \(v-u\in W_0^{1,p}(\Omega)\)。以 \(v-u\) 测试两式之差，得到
\[
 \int_\Omega
 \bigl(|\nabla v|^{p-2}\nabla v-|\nabla u|^{p-2}\nabla u\bigr)
 \cdot(\nabla v-\nabla u)\,\mathrm dx=0.
\]
\cref{b2:lem:H2-monotonicity}说明被积函数非负，且只有两个梯度相同时才为零。因此 \(\nabla(v-u)=0\) 几乎处处，零迹 Poincaré 不等式给出 \(v=u\)。唯一性由此覆盖全部弱解，而不以径向对称为附加条件。
\end{solution}

% 自拟：先在弱形式中缩放，再区分函数源和自然对偶源。
\Needspace{6\baselineskip}
\begin{exercise}\label{b2:ex:H2-scaling}
设 \(d\ge1\)、\(1<p<\infty\)，\(\Omega\subset\mathbb R^d\) 为开集，\(f\in L^1_{\mathrm{loc}}(\Omega)\)，且 \(u\in W^{1,p}_{\mathrm{loc}}(\Omega)\) 满足
\[
 \int_\Omega |\nabla u|^{p-2}\nabla u\cdot\nabla\varphi\,\mathrm dy
 =\int_\Omega f\varphi\,\mathrm dy
 \quad(\varphi\in C_c^\infty(\Omega)).
\]
给定 \(x_0\in\mathbb R^d\)、\(a,r>0\)，令
\(\Omega_r=(\Omega-x_0)/r\)、\(v(x)=a\cdot u(x_0+rx)\)。求 \(v\) 的方程右端，并写出可测集 \(E\subset\Omega_r\) 上梯度能量的换元公式。若 \(\Omega\) 有界、\(u\in W^{1,p}(\Omega)\)，原右端改为 \(\ell\in(W_0^{1,p}(\Omega))^*\) 且弱式对全部零迹测试成立，缩放后的右端泛函应如何定义？
\end{exercise}
\begin{solution}
线性换元的 Sobolev 链式法则给出
\(\nabla v(x)=ar\cdot(\nabla u)(x_0+rx)\)。取 \(\varphi\in C_c^\infty(\Omega_r)\)，令 \(\psi(y)=\varphi((y-x_0)/r)\)。于是 \(\psi\in C_c^\infty(\Omega)\)，且 \(\nabla\psi(y)=r^{-1}(\nabla\varphi)((y-x_0)/r)\)。逐项换元得
\begin{align*}
 \int_{\Omega_r}|\nabla v|^{p-2}\nabla v\cdot\nabla\varphi\,\mathrm dx
 &=a^{p-1}r^{p-d}
   \int_\Omega|\nabla u|^{p-2}\nabla u\cdot\nabla\psi\,\mathrm dy\\
 &=a^{p-1}r^{p-d}\int_\Omega f\psi\,\mathrm dy\\
 &=\int_{\Omega_r}a^{p-1}r^p\cdot f(x_0+rx)\cdot\varphi(x)\,\mathrm dx.
\end{align*}
故新方程为
\[
 -\Delta_pv=a^{p-1}r^p\cdot f(x_0+rx)
 \quad\text{于 }\Omega_r.
\]
右端仍局部可积，计算没有使用二阶导数。梯度能量满足
\[
 \int_E|\nabla v|^p\,\mathrm dx
 =a^pr^{p-d}\int_{x_0+rE}|\nabla u|^p\,\mathrm dy.
\]
任意可测 \(E\) 均可使用此非负积分公式；\(E\) 相对紧时两端有限。

对自然对偶中的源项，不能写成点态复合 \(\ell(x_0+rx)\)。正确的泛函是
\[
 \widetilde\ell(\varphi)
 =a^{p-1}r^{p-d}\cdot
  \ell\bigl(\varphi((\cdot-x_0)/r)\bigr),
 \quad \varphi\in W_0^{1,p}(\Omega_r).
\]
由光滑紧支撑函数的逼近，测试函数的换元把 \(W_0^{1,p}(\Omega_r)\) 映到 \(W_0^{1,p}(\Omega)\)，并有
\[
 \bigl\|\nabla[\varphi((\cdot-x_0)/r)]\bigr\|_{L^p(\Omega)}
 =r^{d/p-1}\|\nabla\varphi\|_{L^p(\Omega_r)}.
\]
这证明 \(\widetilde\ell\) 连续，前面的计算也在零迹测试类中成立。源项随幅度按 \(a^{p-1}\) 变化，能量则按 \(a^p\) 变化；两者分别对应通量和它的凸势。
\end{solution}

% 自拟：可测加权一维通量的积分常数，由端点条件唯一确定。
\Needspace{6\baselineskip}
\begin{exercise}\label{b2:ex:H2-ordered-source}
设 \(1<p<\infty\)，\(a:(0,1)\to\mathbb R\) 可测，且
\(0<a_0\le a(x)\le a_1<\infty\) 几乎处处。构造
\[
 -(a(x)|u'|^{p-2}u')'=1,
 \quad u\in W_0^{1,p}(0,1)
\]
的弱解，把未知积分常数化为一个严格单调的标量方程，并证明解唯一。若 \(a(x)=\kappa>0\) 为常数，求闭式并检查解的非负性。
\end{exercise}
\begin{solution}
记通量为 \(q=a|u'|^{p-2}u'\)。弱方程给出分布等式 \(q'=-1\)，所以 \(q+x\) 在区间上为常数，即 \(q(x)=C-x\)。由于 \(a\ge a_0>0\)，可以逐点反解为
\[
 u'(x)=\operatorname{sgn}(C-x)
       \left|\frac{C-x}{a(x)}\right|^{1/(p-1)}.
\]
取左端点值为零后，只需确定右端点值。令
\[
 H(C)=\int_0^1\operatorname{sgn}(C-t)
                   \left|\frac{C-t}{a(t)}\right|^{1/(p-1)}\,\mathrm dt.
\]
在任意有界的 \(C\) 区间上，被积函数由一个常数控制；它又逐点连续，故控制收敛说明 \(H\) 连续。若 \(C_1<C_2\)，函数 \(s\mapsto\operatorname{sgn}(s)|s|^{1/(p-1)}\) 的严格递增性和 \(a(t)>0\) 说明两个被积函数几乎处处严格有序，因此 \(H(C_1)<H(C_2)\)。同时 \(H(0)<0<H(1)\)，所以恰有一个 \(C\in(0,1)\) 满足 \(H(C)=0\)。

用此常数定义
\[
 u(x)=\int_0^x\operatorname{sgn}(C-t)
                   \left|\frac{C-t}{a(t)}\right|^{1/(p-1)}\,\mathrm dt.
\]
其导数有界，端点值都为零，故 \(u\in W_0^{1,p}(0,1)\)。通量等于 \(C-x\)，任取 \(\varphi\in W_0^{1,p}(0,1)\)，利用绝对连续代表及端点值可得
\[
 \int_0^1a|u'|^{p-2}u'\varphi'\,\mathrm dx
 =\int_0^1(C-x)\varphi'\,\mathrm dx
 =\int_0^1\varphi\,\mathrm dx.
\]
这验证了弱形式。反过来，任何弱解的通量均为 \(C-x\)，而零边界条件都要求 \(H(C)=0\)，所以解唯一。

在 \((0,C)\) 上 \(u'>0\)，在 \((C,1)\) 上 \(u'<0\)，并且两个端点值为零，因此 \(u\ge0\)。常系数时，被积函数关于 \(t=1/2\) 反对称，故唯一的常数是 \(C=1/2\)。直接积分得到
\[
 u(x)=\kappa^{-1/(p-1)}\frac{p-1}{p}
       \left[2^{-p/(p-1)}-|x-1/2|^{p/(p-1)}\right].
\]
括号内非负，与前面的单调性判断一致。
\end{solution}

% 自拟：闭凸障碍类的直接法和单侧变分，不借用光滑边界。
\Needspace{6\baselineskip}
\begin{exercise}\label{b2:ex:H2-obstacle}
设 \(d\ge1\)、\(1<p<\infty\)，\(\Omega\subset\mathbb R^d\) 是有界非空开集，\(X=W_0^{1,p}(\Omega)\) 取梯度范数，\(\ell\in X^*\)。给定可测函数 \(\psi:\Omega\to\mathbb R\)，假设
\[
 K=\{\,v\in X\mid v\ge\psi\ \text{几乎处处}\,\}\ne\varnothing.
\]
在 \(K\) 上极小化
\[
 J(v)=\frac1p\int_\Omega|\nabla v|^p\,\mathrm dx-\ell(v)
\]
是一个\term[ZhangaiWenti]{障碍问题}[obstacle problem]。证明它有唯一极小元 \(u\)，并证明 \(u\in K\) 为极小元的充要条件是
\[
 \int_\Omega|\nabla u|^{p-2}\nabla u\cdot\nabla(v-u)\,\mathrm dx
 \ge\ell(v-u)\quad(v\in K).
\]
证明中不得添加边界光滑假设。
\end{exercise}
\begin{solution}
集合 \(K\) 凸。若 \(v_n\in K\) 且 \(v_n\to v\) 于 \(X\)，零边界 Poincaré 不等式给出 \(L^p\) 强收敛。选取几乎处处收敛的子列，便有 \(v\ge\psi\) 几乎处处，所以 \(K\) 强闭。闭凸集又是弱闭集：位于其外的点可用连续线性泛函分离，而该泛函在弱拓扑中仍连续。这一步只用函数空间的凸性结构，不要求为 \(\partial\Omega\) 定义逐点法向或迹。

由对偶范数的定义，
\[
 J(v)\ge\frac1p\|\nabla v\|_p^p
           -\|\ell\|_{X^*}\cdot\|\nabla v\|_p.
\]
因为 \(p>1\)，右端随梯度范数趋于无穷而趋于无穷，并有有限下界。选取一个 \(v_0\in K\)，则 \(\inf_KJ\le J(v_0)<\infty\)，从而存在有界极小序列。空间 \(X\) 自反，故可抽取弱收敛子列，其极限 \(u\) 由弱闭性仍在 \(K\) 中。梯度的 \(L^p\) 范数弱下半连续，\(\ell\) 弱连续，故 \(J(u)\le\liminf J(v_n)=\inf_KJ\)。若两个不同的函数同为极小元，零迹 Poincaré 不等式保证其梯度不可能几乎处处相同；\(|\xi|^p\) 的严格凸性就使中点的能量严格更小，矛盾。因此极小元唯一。

固定 \(v\in K\)，令 \(h=v-u\)。对 \(0\le t\le1\)，凸性保证 \(u+th\in K\)，于是 \(J(u+th)\ge J(u)\)。差商中积分项的导数由
\[
 C_p\bigl(|\nabla u|+|\nabla h|\bigr)^{p-1}\cdot|\nabla h|
\]
控制；该函数由 Hölder 不等式可积。因此可以把右导数移入积分，得到
\[
 0\le\left.\frac{\mathrm d}{\mathrm dt}J(u+th)\right|_{t=0+}
 =\int_\Omega|\nabla u|^{p-2}\nabla u\cdot\nabla h\,\mathrm dx-\ell(h),
\]
即题中的不等式。

反之，若 \(u\in K\) 满足该不等式，凸函数的切平面不等式给出
\[
 \frac{|\nabla v|^p-|\nabla u|^p}{p}
 \ge |\nabla u|^{p-2}\nabla u\cdot(\nabla v-\nabla u).
\]
积分并减去 \(\ell(v-u)\)，便得 \(J(v)-J(u)\ge0\) 对所有 \(v\in K\) 成立。可行变分只保证 \(t\ge0\) 的方向，故得到单侧不等式；一般不能再用负方向把它变为等式。
\end{solution}

% 自拟：点值泛函的连续性与精确范数，一维点源的通量跳跃。
\Needspace{6\baselineskip}
\begin{exercise}\label{b2:ex:H2-point-source}
设 \(1<p<\infty\)，\(X=W_0^{1,p}(-1,1)\) 取范数 \(\|v'\|_p\)。证明点值泛函 \(\delta_0(v)=v(0)\) 属于 \(X^*\)，并求其范数。求解
\[
 -(|u'|^{p-2}u')'=\delta_0,
 \quad u\in X,
\]
直接验证弱形式，并证明所得函数是
\(J(v)=p^{-1}\int_{-1}^1|v'|^p\,\mathrm dx-v(0)\)
在 \(X\) 上的唯一极小元。
\end{exercise}
\begin{solution}
一维 \(W^{1,p}\) 函数有唯一的绝对连续代表，以下点值均指此代表。对 \(v\in X\)，端点值为零，因而
\[
 |v(0)|=\left|\int_{-1}^0v'\,\mathrm dx\right|
       =\left|-\int_0^1v'\,\mathrm dx\right|.
\]
在两个长度为一的半区间分别使用 Hölder 不等式，得到
\[
 |v(0)|^p\le\int_{-1}^0|v'|^p\,\mathrm dx,
 \quad
 |v(0)|^p\le\int_0^1|v'|^p\,\mathrm dx.
\]
相加可知 \(|v(0)|\le2^{-1/p}\|v'\|_p\)。函数 \(v(x)=1-|x|\) 取等，故 \(\|\delta_0\|_{X^*}=2^{-1/p}\)。点源虽不是普通可积函数，这里仍是自然能量空间上的连续输入。

令 \(c=2^{-1/(p-1)}\)，取 \(u(x)=c(1-|x|)\)。该函数属于 \(X\)，通量在 \((-1,0)\) 上为 \(c^{p-1}=1/2\)，在 \((0,1)\) 上为 \(-1/2\)。任取 \(\varphi\in X\)，由绝对连续性有
\begin{align*}
 \int_{-1}^1|u'|^{p-2}u'\varphi'\,\mathrm dx
 &=\frac12\bigl(\varphi(0)-\varphi(-1)\bigr)
   -\frac12\bigl(\varphi(1)-\varphi(0)\bigr)\\
 &=\varphi(0).
\end{align*}
这正是方程的弱形式，通量向下跳跃一的大小给出右端的单位点源。

对任意 \(v\in X\)，凸性给出
\[
 J(v)-J(u)\ge
 \int_{-1}^1|u'|^{p-2}u'(v'-u')\,\mathrm dx-(v(0)-u(0))=0.
\]
最后一个等号是以 \(v-u\) 测试所得的弱等式，因而 \(u\) 是极小元。严格凸性说明另一极小元必须具有相同导数，零端点值再给出函数相等。特别地，能量最小性与分布中的通量跳跃在同一个自然空间内一致。
\end{solution}

% 自拟：全空间 L^p 假设与能量截断，不调用非线性 Harnack。
\Needspace{6\baselineskip}
\begin{exercise}\label{b2:ex:H2-liouville}
设 \(d\ge1\)、\(1<p<\infty\)，实值函数
\(u\in W^{1,p}_{\mathrm{loc}}(\mathbb R^d)\cap L^p(\mathbb R^d)\)
满足 \(\Delta_pu=0\) 的弱形式。证明 \(u=0\) 几乎处处。本题不使用非线性 Harnack 定理。
\end{exercise}
\begin{solution}
对 \(R>0\)，取 \(\eta_R\in C_c^\infty(B_{2R})\)，使
\(0\le\eta_R\le1\)、\(\eta_R=1\) 于 \(B_R\)，且
\(|\nabla\eta_R|\le C/R\)。函数 \(\eta_R^p\cdot u\) 属于 \(W_0^{1,p}(B_{2R})\)，其支集紧含于该球。由于通量属于 \(L^{p'}(B_{2R})\)，用光滑紧支撑函数在 \(W^{1,p}\) 中逼近，便可把它代入弱式。这里不要求 \(p\) 为整数：函数 \(t\mapsto t^p\) 在 \([0,1]\) 上具有有界导数，链式法则已经足够。

展开测试函数的梯度并移项，得到
\[
 \int\eta_R^p|\nabla u|^p
 =-p\int\eta_R^{p-1}\cdot u\cdot
       |\nabla u|^{p-2}\nabla u\cdot\nabla\eta_R.
\]
右端的绝对值至多为
\[
 p\int(\eta_R|\nabla u|)^{p-1}\cdot|u|\cdot|\nabla\eta_R|
 \le\frac12\int\eta_R^p|\nabla u|^p
       +C_p\int|u|^p\cdot|\nabla\eta_R|^p.
\]
这里使用共轭指数 \(p'=p/(p-1)\)、\(p\) 的带参数 Young 不等式。吸收第一项，并利用截断梯度只出现在环带内，得到
\[
 \int_{B_R}|\nabla u|^p\,\mathrm dx
 \le\frac{C_p}{R^p}\int_{B_{2R}\setminus B_R}|u|^p\,\mathrm dx
 \le\frac{C_p}{R^p}\|u\|_{L^p(\mathbb R^d)}^p.
\]
固定 \(\rho>0\)，令 \(R\ge\rho\) 趋于无穷，可知 \(\int_{B_\rho}|\nabla u|^p=0\)。因此梯度在全空间几乎处处为零。球上的 Poincaré 不等式说明 \(u\) 在每个球中为常数，重叠球又使这些常数一致。全空间的测度无穷，而 \(u\in L^p(\mathbb R^d)\)，故唯一可能的常数是零。
\end{solution}

% 自拟：线性化计算之外加入明确的 p=4 非叠加模型。
\Needspace{6\baselineskip}
\begin{exercise}\label{b2:ex:H2-linearization}
设 \(d\ge2\)、\(1<p<\infty\)，令 \(F(\xi)=|\xi|^{p-2}\xi\)。对 \(\xi\ne0\)，求 \(DF(\xi)\) 的切向、径向特征值，并说明当 \(p\ne2\) 时，特征值比有界为何仍不能保证对所有非零梯度具有固定的一致椭圆常数。

\ProseParagraphBreak

再令 \(p=4\)、\(d=2\)，在环域 \(\Omega=\{\,x\in\mathbb R^2\mid1/2<|x|<2\,\}\) 上考虑
\(u(x)=|x|^{2/3}\)、\(v(x)=x_1\)。证明二者均为弱 \(4\)-调和函数，而 \(u+v\) 不是。
\end{exercise}
\begin{solution}
对分量逐一求导得
\[
 DF(\xi)=|\xi|^{p-2}I+(p-2)|\xi|^{p-4}\xi\otimes\xi.
\]
若 \(h\perp\xi\)，则第二项作用于 \(h\) 为零，故切向特征值是 \(|\xi|^{p-2}\)；在 \(\xi\) 的方向上，两项相加给出径向特征值 \((p-1)|\xi|^{p-2}\)。最大、最小特征值之比为
\[
 \frac{\max(1,p-1)}{\min(1,p-1)}.
\]
它只依赖 \(p\)，但线性一致椭圆性还要求共同幅度有固定的正下界与有限上界。当 \(p>2\) 时，共同因子 \(|\xi|^{p-2}\) 在 \(\xi\to0\) 时趋零，在 \(|\xi|\to\infty\) 时趋于无穷；当 \(1<p<2\) 时两个极限方向相反。因此椭圆比本身不足以给出全梯度范围的统一常数。只有 \(p=2\) 时 \(DF=I\) 恒定。

对于第二问，径向散度公式给出
\[
 r\,|u'(r)|^2u'(r)
 =r\left(\frac23r^{-1/3}\right)^3=\frac8{27},
\]
故 \(\Delta_4u=0\)。函数 \(v=x_1\) 的梯度恒定，也满足 \(\Delta_4v=0\)。两者在环域上光滑，因而经典等式同时给出弱等式。

令 \(w=u+v\)。对光滑函数，直接展开散度有
\[
 \Delta_4w=|\nabla w|^2\Delta w
                 +2\,D^2w[\nabla w,\nabla w].
\]
在 \(x=(0,1)\) 处，径向函数的导数给出
\[
 \nabla w=(1,2/3),\quad
 D^2w=\begin{pmatrix}2/3&0\\0&-2/9\end{pmatrix},\quad
 \Delta w=4/9.
\]
于是
\[
 \Delta_4w(0,1)
 =\frac{13}{9}\cdot\frac49
   +2\left(\frac23-\frac29\cdot\frac49\right)
 =\frac{16}{9}\ne0.
\]
由于该散度连续，它在一个小邻域内保持正值，用非负且非零的紧支撑测试即可否定齐次弱等式。非线性通量确实破坏了调和函数的叠加性，即使所有参与的函数都光滑。
\end{solution}

% 自拟：定量单调性的源差估计，以及幅度缩放对指数的限制。
\Needspace{6\baselineskip}
\begin{exercise}\label{b2:ex:H2-stability}
设 \(d\ge1\)、\(2\le p<\infty\)，\(\Omega\subset\mathbb R^d\) 是有界非空开集，\(X=W_0^{1,p}(\Omega)\) 取梯度范数，\(f,g\in X^*\)。令 \(u_f,u_g\in X\) 分别为 \(-\Delta_pu=f,g\) 的弱解。证明存在只依赖 \(p\) 的常数 \(C_p\)，使
\[
 \|\nabla(u_f-u_g)\|_p
 \le C_p\|f-g\|_{X^*}^{1/(p-1)}.
\]
再证明：在包含零源项的任意邻域中，不能把右端幂指数 \(1/(p-1)\) 提高为更大的数，同时保持常数与 \(f,g\) 无关。
\end{exercise}
\begin{solution}
本题的存在唯一性由\cref{b2:thm:H2-variational}取权重为一、边界数据为零得到。令 \(w=u_f-u_g\in X\)，两弱式相减并以 \(w\) 测试。由\cref{b2:lem:H2-monotonicity}中 \(p\ge2\) 的定量下界，有
\begin{align*}
 c_p\|\nabla w\|_p^p
 &\le\int_\Omega
   \bigl(|\nabla u_f|^{p-2}\nabla u_f
          -|\nabla u_g|^{p-2}\nabla u_g\bigr)
       \cdot\nabla w\,\mathrm dx\\
 &=\langle f-g,w\rangle
 \le\|f-g\|_{X^*}\cdot\|\nabla w\|_p.
\end{align*}
若 \(\|\nabla w\|_p=0\)，结论成立；否则约去一次幂，再取 \((p-1)\) 次方根即可，常数可取 \(C_p=c_p^{-1/(p-1)}\)。

为检验幂指数，取一个非零 \(z\in C_c^\infty(\Omega)\)，并定义
\[
 f_0(\varphi)=\int_\Omega|\nabla z|^{p-2}\nabla z\cdot\nabla\varphi\,\mathrm dx
 \quad(\varphi\in X).
\]
Hölder 不等式说明 \(f_0\in X^*\)，且
\(f_0(z)=\|\nabla z\|_p^p>0\)，所以 \(f_0\ne0\)。对任意 \(t>0\)，通量的齐次性和唯一性给出
\[
 u_{tf_0}=t^{1/(p-1)}z,\quad u_0=0.
\]
假如在零附近存在指数 \(\beta>1/(p-1)\) 和统一常数 \(C\)，使题中估计的右端可换成 \(C\|f-g\|_{X^*}^\beta\)，那么取 \(f=tf_0\)、\(g=0\)，对充分小的 \(t\) 有
\[
 t^{1/(p-1)}\|\nabla z\|_p
 \le C t^\beta\|f_0\|_{X^*}^\beta.
\]
约去左侧的 \(t\) 次幂并令 \(t\downarrow0\)，便迫使 \(\|\nabla z\|_p=0\)，与所选 \(z\) 矛盾。因此幂指数不能提高。特别地，\(p=2\) 时得到线性 Lipschitz 稳定性；\(p>2\) 时，零源项附近的解算子一般只能具有这里的 Hölder 幅度。
\end{solution}
